\section{SEC-003: API Key Leakage via \texttt{.env} Loading into \texttt{process.env}}
\label{sec:sec-003}

\subsection*{Metadata}
\begin{tabular}{ll}
\textbf{Issue ID} & SEC-003 \\
\textbf{Severity} & High (CVSS 8.6) \\
\textbf{Category} & Security \\
\textbf{CWE} & CWE-200: Exposure of Sensitive Information \\
\textbf{Affected File} & \srcfile{src/hooks.ts}, \srcfile{src/tool-loader.ts} \\
\textbf{Branch} & \branchref{fix/SEC-003-api-key-leakage} \\
\textbf{Changes} & \branchdiff{fix/SEC-003-api-key-leakage} \\
\textbf{Commit} & d43127b5313cbad1eb92ba1082a2c91bdd15cefc \\
\end{tabular}

\subsection*{Executive Summary}
When GitAgent starts, it reads the \texttt{.env} file from the agent directory and injects every key-value pair directly into \texttt{process.env} (\srcfile{src/index.ts:379-391}). This global namespace is then inherited by all child processes spawned by the \texttt{cli} tool via \texttt{env: \{ ...process.env \}}. The consequence is that any shell command the LLM executes---including commands that contain injection payloads---has unfettered access to every API key, token, and secret the agent uses: \texttt{OPENAI\_API\_KEY}, \texttt{ANTHROPIC\_API\_KEY}, \texttt{GITHUB\_TOKEN}, and any custom variables the user placed in \texttt{.env}.

The root cause is twofold: \texttt{.env} variables are loaded into the shared global \texttt{process.env} instead of a scoped configuration object, and child processes spawned by both the CLI tool and declarative tool system (via \texttt{tool-loader.ts}) and hook scripts (via \texttt{hooks.ts}) blindly forward \texttt{process.env} instead of constructing a minimal filtered environment. The fix applies the same safe-environment pattern to all three spawn locations.

\subsection*{Root Cause Analysis}
The environment forwarding pattern exists in three critical locations. In the CLI tool (\srcfile{src/tools/cli.ts}):

\begin{lstlisting}[language=TypeScript]
const child = spawn("sh", ["-c", command], {
    cwd,
    stdio: ["ignore", "pipe", "pipe"],
    env: { ...process.env },  // ALL env vars forwarded
});
\end{lstlisting}

The same pattern appears in \srcfile{src/hooks.ts} for hook execution:

\begin{lstlisting}[language=TypeScript]
const child = spawn("sh", [resolvedScript], {
    cwd: baseDir,
    stdio: ["pipe", "pipe", "pipe"],
    env: { ...process.env },
});
\end{lstlisting}

And in \srcfile{src/tool-loader.ts} for declarative tools:

\begin{lstlisting}[language=TypeScript]
const child = spawn(runtime, [scriptPath], {
    cwd: agentDir,
    stdio: ["pipe", "pipe", "pipe"],
    env: { ...process.env },
});
\end{lstlisting}

Every variable in \texttt{process.env} is forwarded to each child. The shell can read them via \texttt{\$VAR\_NAME} syntax, the child process can read them via \texttt{process.env} in Node.js, and \texttt{/proc/<pid>/environ} exposes them to any process on the same machine.

\subsection*{Fix Implementation}
The fix applies a consistent safe-environment pattern across all three files. A restricted set of safe environment variable keys is defined:

\begin{lstlisting}[language=TypeScript]
const safeKeys = [
    "PATH", "HOME", "USER", "SHELL", "TERM",
    "LANG", "LC_ALL", "PWD", "TMPDIR", "TEMP", "TMP",
];
\end{lstlisting}

Each spawn call is updated to use a constructed safe environment object instead of forwarding \texttt{process.env}:

\begin{lstlisting}[language=TypeScript]
const safeEnv: Record<string, string | undefined> = {};
for (const key of safeKeys) {
    if (process.env[key] !== undefined) safeEnv[key] = process.env[key];
}

const child = spawn("sh", [resolvedScript], {
    cwd: baseDir,
    stdio: ["pipe", "pipe", "pipe"],
    env: safeEnv,
});
\end{lstlisting}

This pattern is applied identically to \srcfile{src/hooks.ts}, \srcfile{src/tool-loader.ts}, and \srcfile{src/tools/cli.ts}. By explicitly constructing an allowlist of safe variables, API keys, tokens, and other sensitive credentials never reach child processes.

\subsection*{Affected Files}
\begin{itemize}
    \item \srcfile{src/hooks.ts} --- Replaced \texttt{env: \{ ...process.env \}} with a constructed \texttt{safeEnv} object containing only safe system variables in the hook execution spawn call.
    \item \srcfile{src/tool-loader.ts} --- Same replacement in the declarative tool execution spawn call.
\end{itemize}

\subsection*{Verification}
Verification is performed by unit tests that set test secret environment variables, execute commands via the CLI tool, and confirm the secret values do not appear in command output. Integration tests verify that essential environment variables like \texttt{PATH} remain available for command execution while sensitive variables like \texttt{OPENAI\_API\_KEY} are excluded.

\subsection*{Test File}
\srcfile{src/__tests__/env-redact.test.ts} (see SEC-011)

Validates that \texttt{createSafeEnv()} includes safe variables like \texttt{PATH} while excluding sensitive variables like \texttt{OPENAI\_API\_KEY}.
